\documentclass[12pt]{article}
\usepackage[utf8]{inputenc}
\usepackage[T1]{fontenc}
\usepackage{hyperref}
\title{Tsetsaut Language Resources for K-12 Curriculum Development}
\author{Compiled by Claude Sonnet 5 (claude-sonnet-5) for Andr\'as Kornai}
\date{}
\begin{document}
\maketitle

\section{Introduction}
Tsetsaut is a dormant (extinct) Athabaskan language of British Columbia and Alaska. Almost everything known comes from Franz Boas's 1894 fieldwork with two Tsetsaut people living among the Nisga'a, published with Pliny Goddard in 1924. The last speaker is reported to have lived into the 1920s or 1930s. No school programme exists, and the evidence available is a small word list; a K-12 curriculum in the language is not realistic from current sources.

\section{Primary Educational Resources}
None found. No primer, textbook or school programme for Tsetsaut turned up in the search.

\section{Linguistic and Reference Materials}
\begin{itemize}
\item Wikipedia overview of the Tsetsaut language: \url{https://en.wikipedia.org/wiki/Tsetsaut_language}
\item Glottolog reference record pointing to the Boas and Goddard (1924) material: \url{https://glottolog.org/resource/reference/id/37437}
\item Bibliography of materials on the language (Yukon Native Language / ydli.org): \url{https://ydli.org/biblios/tsetbib.htm}
\item Rosetta Project page for Tsetsaut: \url{https://rosettapanglossia.longnow.org/index.php?title=Txc}
\end{itemize}

\section{Religious and Community Resources}
The Canadian Encyclopedia entry on the Tsetsaut people notes that about 30 people identified as Tsetsaut in British Columbia as of 2019: \url{https://thecanadianencyclopedia.ca/en/article/tsetsaut}. No community language organisation was found.

\section{Digital and Audio Resources}
No audio or digital learning tools were found. The word list itself was not reproduced in the search results.

\section{Curriculum Development Context}
Any work should be led by Tsetsaut descendants and neighbouring Nisga'a and Tahltan communities, who should decide whether and how the historical word list is used. Related living Athabaskan languages may offer more teachable material. Little else can responsibly be said.

\end{document}